\documentclass[12pt]{article}
\usepackage[T1]{fontenc}
\usepackage{lmodern}

\usepackage[margin=1in]{geometry}
\usepackage{amsmath,amssymb,amsthm,mathtools}
\usepackage{booktabs,array,enumitem,microtype}
\usepackage[colorlinks=true,allcolors=blue]{hyperref}

\newtheorem{theorem}{Theorem}[section]
\newtheorem{proposition}[theorem]{Proposition}
\newtheorem{corollary}[theorem]{Corollary}
\newtheorem{lemma}[theorem]{Lemma}
\theoremstyle{definition}
\newtheorem{definition}[theorem]{Definition}
\newtheorem{example}[theorem]{Example}
\theoremstyle{remark}
\newtheorem{remark}[theorem]{Remark}

\newcommand{\ZZ}{\mathbb Z}
\newcommand{\ii}{\mathrm i}
\newcommand{\vol}{\operatorname{Vol}}
\newcommand{\curl}{\operatorname{curl}}
\newcommand{\dist}{\operatorname{dist}}
\newcommand{\Span}{\operatorname{span}}

\title{Coexact Spectral Hydrodynamics on $L(137,1)$:\\
Exact Modes, Triad Couplings, and Decay}
\author{Luqman Mahmood\thanks{Copyright 2026 Luqman Mahmood. This paper is licensed under Creative Commons Attribution 4.0 International (CC BY 4.0); see \url{https://creativecommons.org/licenses/by/4.0/}. Software is excluded from this grant.}}
\date{30 September 2026}
\hypersetup{pdftitle={Coexact Spectral Hydrodynamics on L(137,1): Exact Modes, Triad Couplings, and Decay},pdfauthor={Luqman Mahmood}}

\begin{document}
\maketitle

\begin{abstract}
We study coexact spectral hydrodynamics on the round lens space $L(137,1)$.
Character formulas for spherical space forms give the signed curl gap,
deck projectors, and invariant multiplicities in a fixed orientation and
normalization.  Specializing $SU(2)_L\times SU(2)_R$ recoupling to
deck-invariant coexact modes then yields an exact
$3j$--$3j$--$9j$ interaction coefficient and a fourth-root phase rule in a
specified ladder-canonical basis.  We give explicit nonzero witnesses and a
finite two-band census, and use the coexact gap in the Hodge--Stokes energy
inequality to prove exponential decay and singleton phase and trajectory
attractors.  These results concern one declared Hodge-viscous model; the
phase is basis-dependent, and the analysis neither selects the order $137$
physically nor proves global three-dimensional regularity.
\end{abstract}

\section{Introduction and scope}\label{sec:introduction}

This paper joins three mathematical layers that are often studied separately:
the curl spectrum of a spherical space form, the interaction algebra of
coexact vector harmonics, and the long-time energy behavior of an unforced
viscous evolution.  The geometry is the
round quotient
\[
 M=L(137,1)=S^3/\ZZ_{137}, \qquad
 \vol(M)=\frac{2\pi^2R_c^3}{137}.
\]
The integer $137$ is a mathematical specialization, not a physically selected
parameter.

The spectral ingredients are specializations of established spherical-space-form
character and lens-space $p$-form formulas
\cite{Bar2019,Lauret2018,IkedaYamamoto1979,Ikeda1980Lens,BenAchour2016,
ChandrasekharKendall1957}.  Likewise, the $S^3$ harmonic and
recoupling framework is classical and appears in explicit vector-harmonic
calculations \cite{HamadaHorata2003,BaigueraEtAl2021}.  Helical and Beltrami
decompositions and triad classifications are also classical
\cite{ConstantinMajda1988,Waleffe1992,Waleffe1993}.  Against that background,
the main contribution here is an exact evaluation of the deck-restricted
coexact coupling tensor on $L(p,q_0)$.  The formula separates three
normalization layers, identifies the complete selection algebra, and fixes a
ladder-canonical phase.  Level $137$ is a fixed worked example, also used in
the companion tensor-curl eta calculation; it is not selected by the model.
For $L(137,1)$ we also exhibit exact witnesses and a
finite tower census.  The spectral gap then gives geometry-specific decay and
attractor corollaries through standard energy and evolutionary-system
arguments \cite{CheskidovFoias2006}.

There are inequivalent viscosity operators on a Riemannian manifold
\cite{ChanCzubakDisconzi2017,SamavakiTuomela2020}; throughout we use the Hodge
viscosity.  The nonlinear three-dimensional weak evolution is therefore
treated as a collection of trajectories rather than a globally single-valued
solution semigroup.  Section~\ref{sec:normalization} fixes geometry and
normalizations, Sections~\ref{sec:triads}--\ref{sec:tower} derive the coupling
and phase formulas and their $L(137,1)$ examples, and
Sections~\ref{sec:decay}--\ref{sec:regularity} study decay, attractors, and the
precise regularity limitation.

\section{Geometry, operators, and normalizations}\label{sec:normalization}

Write a unit quaternion as
$q=x_0+x_1\ii+x_2\mathrm j+x_3\mathrm k$ and set
$x=(x_0,x_1,x_2,x_3)$.  We use the ordered imaginary basis
$(e_1,e_2,e_3)=(\ii,\mathrm j,\mathrm k)$, the induced boundary orientation
of $S^3\subset\mathbb R^4$, and $\epsilon_{123}=+1$.  The round
radius-$R_c$ metric on this unit-quaternion model is $g_{R_c}=R_c^2g_1$.
The pointwise-unit left-invariant fields
\[
 v_a(q)=R_c^{-1}q e_a
\]
then satisfy
\[
 \curl v_a=-\frac{2}{R_c}v_a, \qquad a=1,2,3.
\]
Introduce complex coordinates
$z_1=x_0+\ii x_1$ and $z_2=x_2+\ii x_3$.  For $p\ge2$ and a unit
$q_0\in(\ZZ/p\ZZ)^\times$, put $\zeta=e^{2\pi\ii/p}$ and fix the free,
orientation-preserving deck generator
\begin{equation}\label{eq:deck-action}
 \gamma_{p,q_0}(z_1,z_2)=(\zeta z_1,\zeta^{q_0}z_2).
\end{equation}
The quotient orientation is the one inherited from $S^3$.  The pointwise-unit
right-invariant cover frame is $w_a(q)=R_c^{-1}e_aq$ and satisfies
$\curl w_a=+2w_a/R_c$.  Our numerical space has $(p,q_0)=(137,1)$, for which
exactly $w_1(q)=R_c^{-1}\ii q$ survives the deck action, with
\[
 \curl w_1=+\frac{2}{R_c}w_1.
\]
All four are coexact, and the Hodge Laplacian on this sector is
$\Delta_H=\curl^2$.  Thus all four have eigenvalue $4/R_c^2$.

\paragraph{Rank-one eta convention.}
In the convention of \cite{MahmoodP1Collection2026},
the projected rank-one curl/odd-signature operator is $A_1=*d$ on coclosed
one-forms (hence on the coexact sector for these lens spaces).  With the
complex quotient orientation in \eqref{eq:deck-action}, its ordinary eta
invariant, not the reduced value $(\eta+h)/2$, is
$\eta_1(L(p,q_0))=-4s(q_0,p)$.  The same convention takes
$V_{n+2J}\otimes V_n$ as the positive branch and, on changing from unit
radius to $R_c$, scales $A_J\mapsto R_c^{-1}A_J$ while leaving eta at zero
unchanged \cite{APS1975,Bar2019}.  We use this citation only to align the sign,
orientation, and radius conventions; all spectral, interaction, and
hydrodynamic results below are derived within the present paper.

For coefficient calculations we now set $R_c=1$.  Radius is restored only in
operator eigenvalues and decay rates.  On complexified fields the $L^2$
inner product is Hermitian,
\[
 \langle u,v\rangle_{L^2(X)}=\int_X g_{\mathbb C}(u,\overline v)\,dV.
\]
Define the lifted raw trilinear by
\[
 T_{\mathrm{lift,raw}}(u_1,u_2,u_3)
 =\int_{S^3}\det[x,u_1,u_2,u_3]\,dS.
\]
Here the determinant is extended $\mathbb C$-multilinearly in its three field
slots; no field slot is conjugated.  If invariant fields $u_i$ descend to
$\bar u_i$ on $M=S^3/\ZZ_p$, define
\[
 T_{\mathrm{quot,raw}}(\bar u_1,\bar u_2,\bar u_3)
 =p^{-1}T_{\mathrm{lift,raw}}(u_1,u_2,u_3).
\]
The cover- and quotient-normalized modes are distinct:
\[
 \widehat u^{S^3}=\frac{u}{\|u\|_{L^2(S^3)}},\qquad
 \widehat{\bar u}^{M}=\frac{\bar u}{\|\bar u\|_{L^2(M)}};
\]
the lift of $\widehat{\bar u}^{M}$ is $\sqrt p\,\widehat u^{S^3}$.
Accordingly,
\begin{align*}
 T_{\mathrm{lift},L^2}
   &=T_{\mathrm{lift,raw}}(\widehat u_1^{S^3},
      \widehat u_2^{S^3},\widehat u_3^{S^3}),\\
 T_{\mathrm{quot},L^2}
   &=T_{\mathrm{quot,raw}}(\widehat{\bar u}_1^M,
      \widehat{\bar u}_2^M,\widehat{\bar u}_3^M)
    =\sqrt p\,T_{\mathrm{lift},L^2}.
\end{align*}
These are the three normalization layers used below.

\begin{proposition}[Gap-frame calibration]\label{prop:gap-tensor}
For the pointwise-unit gap frame at $p=137$,
\begin{align*}
 T_{\mathrm{lift,raw}}(v_a,v_b,v_c)&=2\pi^2\epsilon_{abc},\\
 T_{\mathrm{quot,raw}}(v_a,v_b,v_c)&=\frac{2\pi^2}{137}\epsilon_{abc},\\
 T_{\mathrm{quot},L^2}(\widehat{\bar v}_a^M,
   \widehat{\bar v}_b^M,\widehat{\bar v}_c^M)
 &=\sqrt{\frac{137}{2\pi^2}}\epsilon_{abc}.
\end{align*}
Every entry containing $w_1$ is zero.  In particular, the gap tensor is
totally alternating rather than a new interaction algebra.
\end{proposition}
\begin{proof}
At $R_c=1$, the ordered ambient determinant satisfies
$\det[x,v_1,v_2,v_3]=1$ pointwise.  Integration over the unit sphere gives
$2\pi^2$, and determinant alternation gives the first displayed identity for
all $a,b,c$.  For a mixed component, the function
$d\mapsto\int_{S^3}\det[x,v_a,v_b,w_d]\,dS$ transforms as a vector under the
adjoint action of left rotations.  Haar invariance makes this vector fixed by
every such rotation; rotations about two independent axes force it to vanish.
Alternation handles the other placements of $w_d$.  Finally, quotient
integration divides the raw value by $137$, while the three quotient-normalized
lifts contribute $137^{3/2}$, giving the stated $L^2$ layer.
\end{proof}
This calibration is consistent with the standard orthonormal $S^3$ frame and
cyclic bracket conventions used in geometric hydrodynamics
\cite{ModinPreston2025}.

\section{Exact coexact spectrum and deck counting}\label{sec:spectrum}

Let $E^{\mathbb R}_{s\ell/R_c}(X)$ denote the real coexact curl eigenspace
\[
 \{u\in C^\infty(TX):\delta u^\flat=0,
       \ \curl u=(s\ell/R_c)u\},
 \qquad \ell\ge2,\quad s\in\{-1,+1\}.
\]
Magnetic and ladder bases below live in its complexification
$E_{s\ell/R_c}(X)=E^{\mathbb R}_{s\ell/R_c}(X)\otimes_{\mathbb R}\mathbb C$.
The deck-fixed bases are stable under complex conjugation, so the complex
dimension of this complexification equals the real multiplicity printed
below.  A mode of type $(\ell,s)$ means a member of this complexified signed
eigenspace.

For $a\ge0$, define the deck-fixed weight count
\[
 F_{137}(a)=\#\{r\in\{-a,-a+2,\ldots,a\}:r\equiv0\pmod{137}\}.
\]
The cyclotomic deck average gives the closed form
\begin{equation}\label{eq:F137}
 F_{137}(a)=
 \begin{cases}
 2\lfloor a/274\rfloor+1,&a\text{ even},\\
 2\lfloor(a+137)/274\rfloor,&a\text{ odd}.
 \end{cases}
\end{equation}

Curl-eigenfield constructions in spherical settings and lens-space quotient
spectral counting are classical
\cite{ChandrasekharKendall1957,IkedaYamamoto1979,Ikeda1980Lens,Bar2019,
Lauret2018}.
On bounded Euclidean balls and concentric spherical shells, Cantarella,
DeTurck, Gluck, and Teytel give a complete curl-eigenfield construction under
the tangent-boundary condition
\cite[Theorems~A--B and Sections~III--V]{CantarellaEtAl2000Curl}.  This is
adjacent spherical-domain literature, not a source for the closed lens-space
multiplicities in Theorem~\ref{thm:spectrum}.

\begin{theorem}[Invariant coexact multiplicities]\label{thm:spectrum}
For a coexact level $\ell\ge2$,
\begin{align}
 \dim E_{+\ell/R_c}&=F_{137}(\ell)(\ell-1),\label{eq:plusdim}\\
 \dim E_{-\ell/R_c}&=F_{137}(\ell-2)(\ell+1),\label{eq:minusdim}
\end{align}
and both signed branches have Hodge--Stokes eigenvalue
$\ell^2/R_c^2$.
\end{theorem}

\begin{proof}
On the cover, the two signed coexact shells are the
$SU(2)_L\times SU(2)_R$ modules
\[
 s=+1:\ (j_L,j_R)=\left(\frac\ell2,\frac{\ell-2}{2}\right),
 \qquad
 s=-1:\ (j_L,j_R)=\left(\frac{\ell-2}{2},\frac\ell2\right).
\]
For $L(137,1)$ the generator~\eqref{eq:deck-action} acts on the left magnetic
factor.  The cyclic character average is therefore the number of doubled
left weights $r=-2j_L,-2j_L+2,\ldots,2j_L$ divisible by $137$.  Counting the
positive and negative multiples of $137$ in this parity lattice gives
\eqref{eq:F137}: an even endpoint contains zero and
$2\lfloor a/274\rfloor$ nonzero weights, whereas an odd endpoint contains
the odd multiples counted by $2\lfloor(a+137)/274\rfloor$.
Multiplication by the untouched right-factor dimensions $\ell-1$ and
$\ell+1$ proves \eqref{eq:plusdim} and \eqref{eq:minusdim}.  Complex
conjugation preserves the fixed subspace, which identifies these complexified
counts with the real dimensions.  Finally, $\Delta_H=\curl^2$ on coexact
fields, giving eigenvalue $\ell^2/R_c^2$.
\end{proof}

The sign labels in Theorem~\ref{thm:spectrum} follow our orientation; reversing
orientation exchanges them.  At $\ell=2$ the split is $1+3$, so the first
 Hodge--Stokes eigenvalue is $4/R_c^2$ with real multiplicity four.  This
 multiplicity is specific to the exact round metric and is not asserted to
 persist under metric perturbation.  For closed oriented three-manifolds,
Kepplinger proves simplicity of the curl spectrum along generic
one-parameter paths of metrics
\cite[Theorem~1.1]{Kepplinger2024}.  Directly
summing \eqref{eq:plusdim}--\eqref{eq:minusdim}, the cumulative real coexact
 dimensions through the absolute-curl cutoffs
 $|\curl|\le136/R_c,137/R_c,138/R_c$ are $9384,9656,9932$.  For even
 $n\le136$ the cumulative count through $|\curl|\le n/R_c$ is
 $n(n+2)/2$.

The spectrum is infinite.  Its leading Weyl coefficient is $1/137$ of the
cover coefficient because the quotient volume is $1/137$ of the cover volume;
this is an asymptotic statement, not equality of every finite shell.  At every
 signed cover shell the exact deck projector
 \[
 P_{p,q_0}^{(\ell,s)}
 =\frac1p\sum_{k=0}^{p-1}
   \rho_{\ell,s}(\gamma_{p,q_0})^k
 \]
has image equal to the deck-invariant subspace.  Here
$\rho_{\ell,s}$ is the unitary pullback action of the generator
\eqref{eq:deck-action} on $E_{s\ell/R_c}(S^3)$.  Coordinates within a
 degenerate higher shell still require additional basis and phase choices.

For $X\in\{S^3,L(137,1)\}$ with unit curvature radius $R_c=1$, let $C_X$
denote the real-scalar, inhomogeneous single-constant Sobolev constant
\[
 C_X:=\sup_{0\ne f\in H^1(X;\mathbb R)}
 \frac{\|f\|_{L^6(X)}}
 {\bigl(\|\nabla f\|_{L^2(X)}^2+\|f\|_{L^2(X)}^2\bigr)^{1/2}}.
\]
All norms use the unnormalized Riemannian volume of the inherited round metric.
Beckner's critical sphere inequality uses normalized surface measure
$d\xi=dV/(2\pi^2)$
\cite[p.~4816 and p.~4818, Eq.~(8)]{Beckner1992Sphere}.  Converting it to
unnormalized volume gives
\[
 \|f\|_{L^6(S^3)}^2
 \le (2\pi^2)^{-2/3}
 \left(\|f\|_{L^2(S^3)}^2+\frac43\|\nabla f\|_{L^2(S^3)}^2\right).
\]
The constant-function test and $4/3>1$ therefore give the explicit bracket
\begin{equation}\label{eq:s3-sobolev-bracket}
 (2\pi^2)^{-1/3}
 \le C_{S^3}
 \le \frac2{\sqrt3}(2\pi^2)^{-1/3}.
\end{equation}
The isometric covering lift, together with the constant-function test on the
quotient of volume $2\pi^2/137$, gives
\begin{equation}\label{eq:l137-sobolev-bracket}
 \left(\frac{137}{2\pi^2}\right)^{1/3}
 \le C_{L(137,1)}
 \le \frac2{\sqrt3}\left(\frac{137}{2\pi^2}\right)^{1/3},
 \qquad
 C_{L(137,1)}\le137^{1/3}C_{S^3}.
\end{equation}
The ratio of the displayed upper and lower endpoints is $2/\sqrt3$ in both
brackets.  These are explicit two-sided bounds only: no endpoint is identified
as the best constant, no equality case is asserted, and no sharpness claim is
made inside either bracket.

\section{Exact deck-restricted coexact triads}\label{sec:triads}

For a type $(\ell,s)$, as defined in Section~\ref{sec:spectrum}, set
\[
 (j_L,j_R)=
 \begin{cases}
 ((\ell-2)/2,\ell/2),&s=-1,\\
 (\ell/2,(\ell-2)/2),&s=+1.
\end{cases}
\]
The coordinate torus acts by
$(z_1,z_2)\mapsto(e^{\ii\alpha}z_1,e^{\ii\beta}z_2)$.
A polynomial vector field has weight $(w_1,w_2)\in\ZZ^2$ when its
pushforward by this torus action transforms by
$e^{\ii(w_1\alpha+w_2\beta)}$.  Pullback weights have the opposite signs.
In the Condon--Shortley magnetic basis,
\[
 w_1=m_L+m_R,\qquad w_2=m_L-m_R,
 \qquad (2m_L,2m_R)=(w_1+w_2,w_1-w_2),
\]
where $m_L=-j_L,-j_L+1,\ldots,j_L$ and likewise for $m_R$.
The weights are integers for the coexact modules above.  For $p\ge2$ and
$q_0\in(\ZZ/p\ZZ)^\times$, the mode descends under
\eqref{eq:deck-action} exactly when
\begin{equation}\label{eq:deck}
 w_1+q_0w_2\equiv0\pmod p.
\end{equation}
All numerical results below specialize this free action to $(p,q_0)=(137,1)$.

Helical/Beltrami decompositions and Fourier-space triad classifications are
classical \cite{ConstantinMajda1988,Waleffe1992,Waleffe1993}.  We specialize
the corresponding recoupling machinery to deck-restricted coexact modes and
derive the resulting exact selection algebra.

Call a type triple \emph{representation-admissible} when the two Wigner rows
and the three spin-one columns obey their triangle and integrality conditions;
in particular $\ell_1+\ell_2+\ell_3$ is even.  We sort types
lexicographically, with $-1<+1$.  Representation admissibility is only a
support condition: the $9j$ or a magnetic $3j$ may still vanish exactly.

\begin{theorem}[Convention-complete coefficient]\label{thm:triad}
Let three ladder-canonical, cover-normalized modes
$\widehat u_i^{S^3}$ form a representation-admissible triple and satisfy
\eqref{eq:deck}.  Define
\begin{equation}\label{eq:lambda}
 \Lambda^2=
 3\prod_{i=1}^3(2j_{L,i}+1)(2j_{R,i}+1)
 \begin{Bmatrix}
 j_{L,1}&j_{L,2}&j_{L,3}\\
 1&1&1\\
 j_{R,1}&j_{R,2}&j_{R,3}
 \end{Bmatrix}^{\!2}.
\end{equation}
If the $9j\,3j\,3j$ product is zero, then
$T_{\mathrm{lift},L^2}=0$ and no phase is assigned.  Otherwise, first sort the
three modes and define $\phi$ by Theorem~\ref{thm:phase}; for that sorted
order,
\begin{align}
 \pi T_{\mathrm{lift},L^2}
   (\widehat u_1^{S^3},\widehat u_2^{S^3},\widehat u_3^{S^3})
 ={}&\phi(\ell_i,s_i)\sqrt{\Lambda^2}
 \begin{pmatrix}
 j_{L,1}&j_{L,2}&j_{L,3}\\
 m_{L,1}&m_{L,2}&m_{L,3}
 \end{pmatrix}\notag\\[-1mm]
 &\times
 \begin{pmatrix}
 j_{R,1}&j_{R,2}&j_{R,3}\\
 m_{R,1}&m_{R,2}&m_{R,3}
 \end{pmatrix}.\label{eq:fulltriad}
\end{align}
An arbitrary ordering is defined by sorting and multiplying by the sign of
the sorting permutation.  Equivalently, the paired $3j$ symbols provide that
odd-permutation sign while the reduced phase is permutation-invariant on the
admissible domain.  Thus the full trilinear is alternating, and its exact
nonvanishing criterion is the nonvanishing of the displayed Wigner product.
\end{theorem}

\begin{proof}
Use Peter--Weyl matrix coefficients normalized in
$L^2(S^3,dS)$ and the real Condon--Shortley spin-one column intertwiner whose
highest nonzero Clebsch--Gordan coefficient is positive.  Contract the three
column indices with the positively oriented alternating invariant fixed by
$\epsilon_{123}=+1$.  The defining $9j$ recoupling identity then gives, in
this real reference, the signed reduced coefficient
\[
 C_{\mathrm{ref}}=
 \sqrt{3\prod_{i=1}^3(2j_{L,i}+1)(2j_{R,i}+1)}
 \begin{Bmatrix}
  j_{L,1}&j_{L,2}&j_{L,3}\\
  1&1&1\\
  j_{R,1}&j_{R,2}&j_{R,3}
 \end{Bmatrix}.
\]
The magnetic contractions left after recoupling are precisely the two
Condon--Shortley $3j$ symbols in~\eqref{eq:fulltriad}.  Converting normalized
Haar matrix coefficients to the $dS$-normalized fields and evaluating the
remaining scalar on the oriented gap frame of
Proposition~\ref{prop:gap-tensor} fixes the global factor to $1/\pi$; this
gives $\pi T=C_{\mathrm{ref}}(3j_L)(3j_R)$.  Squaring the real reduced
magnitude yields~\eqref{eq:lambda}, while Theorem~\ref{thm:phase} restores
the exact ray and $9j$ signs.  On the admissible domain
$J_L+J_R=\sum_i\ell_i-3$ is odd.  Hence the two $3j$ symbols together change
by the sign of an odd permutation, whereas $\phi$ does not, proving
alternation.  Deck averaging acts on each mode separately and gives
\eqref{eq:deck}.  This is the standard $S^3$ recoupling mechanism
\cite{HamadaHorata2003,BaigueraEtAl2021} with every normalization and row
order fixed.
\end{proof}

Deck validity, triangle conditions, parity, and magnetic-weight sums are only
necessary filters.  For example,
\[
 \begin{pmatrix}2&3&3\\0&-2&2\end{pmatrix}=0
\]
despite coarse support.  Thus ``supported'' never means ``nonzero.''

\section{The ladder-canonical phase}\label{sec:phase}

\begin{lemma}[Adjacent-row $9j$ sign]\label{lem:9j-sign}
Let the top row be $(a,b,c)$ and the bottom row be
$(g,h,i)=(a+\sigma_1,b+\sigma_2,c+\sigma_3)$ with
$\sigma_r\in\{-1,+1\}$.  In a mixed shift pattern place the unique unlike
shift in the middle column before taking the Racah expansion below.
If $k$ is the number of plus shifts, then every
surviving term in the Racah expansion has sign $+1$ for $k=0,3$ and sign
$-1$ for $k=1,2$.  Consequently the special $9j$ symbol has that sign when
its Racah support is nonempty and is structurally zero when the support is
empty.  The statement includes boundary rows, where one of the candidate
summands may disappear.
\end{lemma}

\begin{proof}
Column permutations have phase $+1$, since the total $9j$ spin sum is
$2(a+b+c)+2k$.  Thus the unique shift of a mixed pattern may be placed in the
middle column.  In the exact Racah expansion
\[
 \sum_x(-1)^{2x}(2x+1)
 \begin{Bmatrix}a&b&c\\1&i&x\end{Bmatrix}
 \begin{Bmatrix}1&1&1\\b&x&h\end{Bmatrix}
 \begin{Bmatrix}g&h&i\\x&a&1\end{Bmatrix},
\]
the middle $6j$ restricts $x=b+d$.  Substitution in the Racah bounds gives
the complete table below.  Here $S=a+b+c$, $G=g+h+i$, and
$(z_1,z_2,z_3)$ are the three Racah upper summation integers:
\[
\begin{array}{c|c|c|c|c}
 (\sigma_1,\sigma_2,\sigma_3)&k&d
 &(z_1-S,z_2-2b,z_3-G)&\text{sign}\\ \hline
 (-,-,-)&0&-1,0&(0,0,1),(0,1,2)&+,+\\
 (+,+,+)&3&0,1&(1,2,0),(2,3,0)&+,+\\
 (-,+,-)&1&0,1&(0,2,0),(0,3,1)&-,-\\
 (+,-,+)&2&-1,0&(0,0,0),(1,1,0)&-,-
\end{array}
\]
up to the harmless column permutations just described.  Every triangle
prefactor, factorial denominator, and $2x+1$ is positive; the displayed signs
are the parity of $2x+z_1+z_2+z_3$ in the three single-summand $6j$ factors.
Hence all surviving terms have one sign and cannot cancel.  At a boundary an
inadmissible value of $d$ simply drops out, leaving the same conclusion.
\end{proof}

\begin{theorem}[All-level phase rule]\label{thm:phase}
For a sorted representation-admissible type triple whose Wigner product in
Theorem~\ref{thm:triad} is nonzero, put
\[
 H=\frac{\ell_1+\ell_2+\ell_3}{2},\qquad
 n_-=\#\{i:s_i=-1\}.
\]
In the specified ladder-canonical convention, every nonzero reduced triad has
\begin{equation}\label{eq:phase}
 \phi=\ii^{\,H+3(n_-\bmod2)}.
\end{equation}
\end{theorem}

\begin{proof}
In ambient coordinates $x=(x_0,x_1,x_2,x_3)$, set
$z=x_0-\ii x_1$ and $a_\pm=x_2\pm\ii x_3$.  The lexicographically normalized
highest rays are
\[
 u_{\ell,+}=\ii^{\ell-2}z^{\ell-2}
    (\ii a_-,a_-,-\ii z,-z),\qquad
 u_{\ell,-}=\ii^{\ell-2}z^{\ell-2}
    (-\ii a_+,-a_+,\ii z,-z).
\]
Here the component order is $x_0,x_1,x_2,x_3$, and monomials in each
component are ordered by descending lexicographic order
$x_3>x_2>x_1>x_0$.  The first nonzero coefficient is then $+1$.
To fix the remaining ladder phases explicitly, write $w_a=e_aq$ and
$v_a=qe_a$ for the ambient quaternion frames of Section~2, with
$(e_1,e_2,e_3)=(\ii,\mathrm j,\mathrm k)$, and let
$\mathcal L_Xu=[X,u]$.  Our lowering operators are
\[
 J_{L,-}=\tfrac12\mathcal L_{w_2+\ii w_3},\qquad
 J_{R,-}=\tfrac{\ii}{2}\mathcal L_{v_2-\ii v_3}.
\]
Starting from the displayed highest rays, normalize each lowered vector by
the positive coefficient $\sqrt{(j+m)(j-m+1)}$.  Thus
``Condon--Shortley'' includes these explicit operator phases, not only the
magnitudes of the ladder coefficients.

For completeness, set $\zeta_8=e^{\ii\pi/4}$ and choose the coherent real
Condon--Shortley reference
\[
 f_n=\zeta_8^n z^n,\qquad
 E_+=-\ii(\ii a_-,a_-,-\ii z,-z),\qquad
 E_-=-(-\ii a_+,-a_+,\ii z,-z).
\]
Maximal scalar products satisfy $f_nf_m=f_{n+m}$ and all lowering
coefficients are positive.  Comparison with the explicit highest rays gives
the per-mode phase
\[
 \eta_{\ell,s}=\zeta_8^{\ell+2\mathbf1_{\{s=-1\}}}.
\]
Thus, with $k=n_-$,
\[
 \prod_i\eta_{\ell_i,s_i}
 =\zeta_8^{\sum_i\ell_i+2k}=\ii^{H+k}.
\]
For these adjacent rows a minus-helicity mode is a plus shift, so
Lemma~\ref{lem:9j-sign} contributes $(-1)$ exactly for $k=1,2$.  Therefore
\[
 \phi=\ii^{H+k}(-1)^{\mathbf1_{\{k=1,2\}}}
      =\ii^{H+3(k\bmod2)},
\]
which is~\eqref{eq:phase}.  If the Racah support or either magnetic $3j$
vanishes, Theorem~\ref{thm:triad} assigns the coefficient zero rather than an
undefined phase.  Finally, on the admissible domain
$J_L+J_R=\sum_i\ell_i-3$ is odd: the reduced phase is invariant under column
permutation, while the paired $3j$ symbols contribute the odd-permutation
sign.  This is the exchange covariance required by the alternating geometric
trilinear.
\end{proof}

Theorem~\ref{thm:phase} is not a physical or basis-free phase statement.
Rephasing individual modes changes individual coefficient phases.  It says
that this fully specified basis has an exact discrete phase law.

\section{Exact \texorpdfstring{$L(137,1)$}{L(137,1)} witnesses and finite tower data}\label{sec:tower}

This section provides exact regression examples for the coefficient formula
and a finite comparison of its two specified frequency bands.  The data are
not an asymptotic cascade model.  The actual deck restriction acts on magnetic
weights, not on the level sum: the additional filter
$\ell_1+\ell_2+\ell_3\equiv0\pmod{137}$ would discard valid triads even
when every $\ell_i\ge137$.  The following explicit triad proves this.  At
$R_c=1$, set $w_\pm=w_2\pm\ii w_3$.  The modes
\[
 u_1=z_1^{135}w_+,\qquad u_2=\overline{z}_1^{135}w_-
\]
have curl $+137$ and weights
$(2m_L,2m_R)=(-137,-135)$ and $(137,135)$ in this pushforward convention.

To fix the scale of the third mode, for the three prefixes used below define
the zonal harmonic
\[
 Z_N[P]=P\sum_{k=0}^Nc_k|z_1|^{2k}|z_2|^{2(N-k)}
\]
as the unique harmonic polynomial of this form normalized by $c_N=1$.  Put
\[
 f_0=Z_{68}[1],\qquad f_+=Z_{67}[z_1z_2],\qquad
 f_-=Z_{67}[\overline z_1\,\overline z_2]
\]
and choose the concrete weight-$(0,0)$ curl-$+138$ representative
\begin{equation}\label{eq:counterexample-u3}
 u_3=-f_+w_-+\frac{\ii}{34}f_0w_1+f_-w_+.
\end{equation}
Thus no free ray scale remains.  All three modes descend to $L(137,1)$, and
exact integration gives
\[
 T_{\mathrm{lift,raw}}(u_1,u_2,u_3)
 =\frac{187232862460210808175}
 {1063298216584182804321625733272656899038}\pi^2\ne0.
\]
Since $137+137+138=412\equiv1\pmod{137}$, the congruence is not necessary.
The exact criterion remains the Wigner product in
Theorem~\ref{thm:triad}; this example makes no separate statement about
sufficiency.

For a direct signed tower anchor take
\[
 ((\ell,s)_1,(\ell,s)_2,(\ell,s)_3)
 =((137,+),(139,-),(274,+)),
\]
with
\[
 2m_L=(137,137,-274),\qquad 2m_R=(135,137,-272).
\]
The raw magnitude below is attached to concrete representatives, rather than to
three freely rescalable rays.  Let
\[
 J_{R,-}=J_{R,1}-\ii J_{R,2},\qquad
 J_{R,-}|j,m\rangle=\sqrt{(j+m)(j-m+1)}\,|j,m-1\rangle
\]
be the Condon--Shortley-normalized right lowering operator, and set
\[
 U_1=u_{137,+},\qquad U_2=J_{R,-}u_{139,-},\qquad
 U_3=\overline{u_{274,+}}.
\]
These representatives have exactly the displayed magnetic weights.  If
$\|U_i\|_{L^2(S^3)}^2=\pi^2N_i$, direct exact integration gives
$N_1N_2N_3=556/320229$.  The lifted, cover-unnormalized representative
integral has magnitude
\[
 \left|T_{\mathrm{lift,raw}}(U_1,U_2,U_3)\right|=\frac{180}{12467}\pi^2,
\]
independently of the phase chosen for the ambient right lowering operator.
The cover-normalized lifted square is
\[
 \pi^2\left|T_{\mathrm{lift},L^2}
 (\widehat U_1^{S^3},\widehat U_2^{S^3},\widehat U_3^{S^3})\right|^2
 =\frac{28503900}{237409081}.
\]
They are related, not identical:
\begin{equation}\label{eq:anchor-conversion}
 \left(\frac{180}{12467}\right)^2
 =\frac{28503900}{237409081}\frac{556}{320229}.
\end{equation}

The finite window
\[
 \{137,138,139,140\}\cup\{274,275,276,277\}
\]
contains both helicities.  A reduced key consists of a lexicographically
sorted type triple together with a doubled-left-weight triple
$c=(2m_{L,1},2m_{L,2},2m_{L,3})$.  Each $c_i$ lies in its spin lattice, is
divisible by $137$, and satisfies $c_1+c_2+c_3=0$.  For repeated adjacent
types we impose $c_1\ge c_2$ or $c_2\ge c_3$, respectively, to quotient only
the corresponding exchange duplication.  Put
\[
 Q(e)=\Lambda(e)^2
 \begin{pmatrix}j_{L,1}&j_{L,2}&j_{L,3}\\
                 m_{L,1}&m_{L,2}&m_{L,3}\end{pmatrix}^{\!2}.
\]
The reduced amplitude is total on every stored key, but no phase is assigned
to a structural zero:
\[
 A_{\mathrm{red}}(e)=
 \begin{cases}
  0,&Q(e)=0,\\
  \displaystyle\phi(e)\sqrt{\Lambda(e)^2}
  \begin{pmatrix}j_{L,1}&j_{L,2}&j_{L,3}\\
                  m_{L,1}&m_{L,2}&m_{L,3}\end{pmatrix},&Q(e)>0,
 \end{cases}
 \qquad
 R(e)=|A_{\mathrm{red}}(e)|^2.
\]
Thus $R(e)=\Lambda(e)^2(3j_L(e))^2$, while the right magnetic dependence
remains factored:
$\pi T_{\mathrm{lift},L^2}=A_{\mathrm{red}}(e)(3j_R)$.

There are $816$ sorted type triples in the window.  Of these, $171$
cross-band types and $58$ all-high-band types have a nonzero reduced key.
Let $\mathcal B_X$ and $\mathcal B_{H2}$ be their respective nonzero reduced
key sets; they contain $385$ and $160$ entries.  The inherited first-band
block $\mathcal B_0$ stores $56$ reduced keys, $46$ with nonzero
$A_{\mathrm{red}}$.  Expanding the right-weight $3j_R$ factor gives
$806{,}408$ eligible indexed completion slots across all $56$ stored keys,
including keys whose reduced factor vanishes.  It is therefore neither a
right-resolved enumeration nor a count of nonzero pointwise coefficients.

For $B\in\{0,X,H2\}$ define the coherent reduced sum
\[
 S_B=\sum_{e\in\mathcal B_B}A_{\mathrm{red}}(e)
     =\sum_{(r,a,b)\in\mathcal L_B}(a+\ii b)\sqrt r.
\]
Equal positive rational radicands are aggregated, and zero aggregate
coefficients are removed.  The exact-arithmetic supplement gives the complete
canonical lists $\mathcal L_0$, $\mathcal L_X$, and $\mathcal L_{H2}$, with
$28$, $149$, and $73$ rows, respectively, together with directed rational
enclosures for every sum.  Numerically,
\begin{align*}
 S_0&\approx-0.215211944909057
       +1.72881536303333\times10^{-7}\ii,\\
 S_X&\approx40.3109087370940+0.444138023269106\ii,\\
 S_{H2}&\approx15.4464538323349-0.0742113043164842\ii.
\end{align*}

Define
\[
 M_{H1}=\max_{e\in\mathcal B_0}R(e),\qquad
 M_{H2}=\max_{e\in\mathcal B_{H2}}R(e).
\]
With the canonical key ordering, the first-band maximizer is
\[
 ((138,-),(140,-),(140,-)),\qquad 2m_L=(0,0,0).
\]
The two sign-related high-band maximizers have
\[
 ((276,-),(277,-),(277,-)),\qquad
 2m_L=(-274,137,137)\ \text{or}\ (274,-137,-137).
\]
The exact statistic $M_{H2}/M_{H1}=N/D$ is recorded in the supplement.  Its
decimal begins $20.508403$, and the corresponding reduced-amplitude ratio is
\[
 \sqrt{M_{H2}/M_{H1}}\approx4.52862051343485.
\]
All tower counts, ratios, and cancellations in this section are finite-window,
convention-dependent data.  They imply no asymptotic growth, monotonicity,
cascade, turbulence, or typicality law.

\section{Unforced Hodge--Stokes decay}\label{sec:decay}

The Leray--Hopf weak framework and energy--Poincar\'e decay mechanism are
classical \cite{Leray1934,Hopf1950,Temam1997,CheskidovFoias2006}.  Here we
specialize them to the declared Hodge operator and its exact quotient gap.

We now fix the viscous model.  Let
\[
 H=\overline{\{v\in C^\infty(TM):\operatorname{div}v=0\}}^{\,L^2},
 \qquad \mathbb P:L^2(TM)\longrightarrow H
\]
be the real solenoidal phase space and Leray projector.  Musical
identifications are understood.  With the positive-sign convention
$\Delta_H=d\delta+\delta d$, define
\[
 A=\mathbb P\Delta_H,\qquad
 D(A)=H^2(TM)\cap H,\qquad V=D(A^{1/2}).
\]
Since $H^1(M;\mathbb R)=0$, this is the positive Hodge--Stokes operator on
the coexact sector.  We study only
\begin{equation}\label{eq:hodge-ns}
 \partial_tu+\mathbb P(\nabla_u u)+\nu Au=0,
 \qquad u(0)=u_0\in H,\qquad \nu>0.
\end{equation}
No conclusion below is transferred to a rough-Laplacian or
deformation-tensor viscosity.

An admissible Leray--Hopf trajectory is a distributional solution
\[
 u\in C_{\mathrm w}([0,\infty);H)
 \cap L^\infty_{\mathrm{loc}}([0,\infty);H)
 \cap L^2_{\mathrm{loc}}([0,\infty);V)
\]
for which there is a full-measure set of admissible starting times
$\mathcal G_u\subset[0,\infty)$, containing the prescribed initial time
$0$, such that
\begin{equation}\label{eq:energy-inequality}
 \|u(t)\|_H^2+2\nu\int_s^t\|A^{1/2}u(r)\|_H^2\,dr
 \le\|u(s)\|_H^2,
 \qquad s\in\mathcal G_u,\quad t\ge s.
\end{equation}
For complete trajectories the same definition is used on $\mathbb R$, with
$\mathcal G_u\subset\mathbb R$ of full measure.
For each datum in $H$, such forward trajectories are obtained by the
standard solenoidal Galerkin construction \cite{Leray1934,Hopf1950,Temam1997}:
project onto the coexact eigenfields, use the cancellation
$\langle\nabla_u u,u\rangle=0$, and pass to a subsequence using compact
$H^1\hookrightarrow L^2$ on the closed quotient. The energy bound and weak
continuity give the stated admissible-time inequality. This supplies
existence, not uniqueness or global strong regularity.

The coexact Poincar\'e inequality is
\[
 \|A^{1/2}v\|_H^2\ge\lambda_1\|v\|_H^2,
 \qquad\lambda_1=4/R_c^2.
\]
Combining this inequality with~\eqref{eq:energy-inequality}, applying the
integral Gronwall argument at admissible starting times, and using weak
continuity and lower semicontinuity at terminal times gives, for every
$s\in\mathcal G_u$ and $t\ge s$,
\begin{equation}\label{eq:energy-decay}
 \|u(t)\|_H^2\le
 \|u(s)\|_H^2e^{-2\nu\lambda_1(t-s)}
 =\|u(s)\|_H^2\exp\!\left[-\frac{8\nu(t-s)}{R_c^2}\right].
\end{equation}
Hence
\begin{equation}\label{eq:norm-decay}
 \|u(t)\|_H\le
 \|u(s)\|_H\exp\!\left[-\frac{4\nu(t-s)}{R_c^2}\right].
\end{equation}

\section{Evolutionary and trajectory attractors}\label{sec:attractors}

We use standard dissipative and evolutionary-system language
\cite{Temam1997,CheskidovFoias2006}, while retaining set-valued dynamics
because uniqueness is unavailable.

Let $\mathcal E$ be the collection of all admissible Leray--Hopf
trajectories.  For bounded $B\subset H$, define the states reachable at time
$t$ by
\[
 \mathcal U(t)B=\{u(t):u\in\mathcal E,\ u(0)\in B\}.
\]
This notation is set-valued; no semigroup or composition law is asserted.
For $X,Y\subset H$, let
\[
 \dist_H(X,Y)=\sup_{x\in X}\inf_{y\in Y}\|x-y\|_H
\]
be the one-sided Hausdorff semidistance.  If
$R_B=\sup_{b\in B}\|b\|_H$, then
\begin{equation}\label{eq:phase-attraction}
 \dist_H(\mathcal U(t)B,\{0\})
 \le R_Be^{-4\nu t/R_c^2}.
\end{equation}

Let $\mathcal K^+$ be the forward trajectories, endowed with the local weak
topology inherited from $C_{\mathrm{loc}}([0,\infty);H_{\mathrm w})$.
Define $(\tau_hu)(r)=u(r+h)$ for $h\ge0$.  Let
$\mathcal K\subset C_{\mathrm{loc}}(\mathbb R;H_{\mathrm w})$ be the complete
trajectories bounded in $H$, with the analogous translation action.

\begin{theorem}[Zero phase and trajectory attractors]\label{thm:attractor}
The strong $H=L^2$ phase attractor is $\{0\}$.  Moreover, $\mathcal K$
consists only of the constant-zero complete trajectory, and its restriction
to $[0,\infty)$ is the trajectory attractor in the local weak topology.
\end{theorem}

\begin{proof}
Equation~\eqref{eq:phase-attraction} proves uniform strong attraction of every
bounded phase set.  Zero initial energy forces the zero trajectory by
\eqref{eq:energy-inequality}, so the singleton is invariant; every closed
attracting set must contain it.

If $u\in\mathcal K$ and
$M=\sup_{r\in\mathbb R}\|u(r)\|_H$, then for fixed $t$ choose
$s_n\in\mathcal G_u$ with $s_n<t$ and $s_n\to-\infty$.  Equation
~\eqref{eq:energy-decay} gives
\[
 \|u(t)\|_H^2\le M^2e^{-8\nu(t-s_n)/R_c^2}\longrightarrow0.
\]
Thus $u(t)=0$.  Finally, for every $T>0$, trajectories whose initial values
have $H$-radius at most $R$ satisfy
\[
 \sup_{0\le r\le T}\|(\tau_tu)(r)\|_H
 \le Re^{-4\nu t/R_c^2}.
\]
The positive restriction of the zero complete trajectory therefore attracts
in $C_{\mathrm{loc}}([0,\infty);H_{\mathrm w})$, in fact with the displayed
strong bound.  Its dimension is zero because it is a singleton, not because
of a spectral cutoff.
\end{proof}

\section{Regularity boundary and conditional continuation}\label{sec:regularity}

The energy-class interpolation is the standard Ladyzhenskaya bound, while the
regularity comparison is the classical Serrin criterion
\cite{Ladyzhenskaya1969,Serrin1962}.  Use the Hodge graph norm
\[
 \|v\|_{H_A^1}^2=\|v\|_H^2+\|A^{1/2}v\|_H^2,
 \qquad v\in V.
\]
At unit curvature radius, the positive-sign Weitzenb\"ock convention is
$\Delta_H=\nabla^*\nabla+\mathrm{Ric}$ with $\mathrm{Ric}=2g$ on one-forms
\cite[Section~2.2, Theorems~2.6--2.7]{Aino2021}.  Hence, for every smooth
coexact one-form $v$,
\[
 \|A^{1/2}v\|_2^2=\|\nabla v\|_2^2+2\|v\|_2^2.
\]
The quadratic-form Kato inequality
\cite[Corollary~2.5]{Guneysu2014Kato} then gives
\[
 \|\nabla|v|\|_2^2+\|v\|_2^2
 \le \|\nabla v\|_2^2+\|v\|_2^2
 \le \|v\|_{H_A^1}^2.
\]
For $v\in V$, choose smooth coexact $v_n\to v$ in the graph norm.  Applying
the closed Weitzenb\"ock form identity to the graph-norm convergence gives
$v_n\to v$ in $L^2$ and
$\nabla v_n\to\nabla v$ in $L^2$.  Kato bounds $|v_n|$ in scalar $H^1$, while
$\bigl||v_n|-|v|\bigr|\le|v_n-v|$ gives strong $L^2$ convergence to $|v|$.
A subsequence converges weakly in $H^1$ to that same limit.  Weak lower
semicontinuity, followed by
$\|\nabla v_n\|_2\to\|\nabla v\|_2$, gives
$\|\nabla|v|\|_2\le\|\nabla v\|_2$.  The scalar Sobolev inequality therefore
passes to the graph-form domain, giving
\begin{equation}\label{eq:graph-embedding}
 C_V:=\frac2{\sqrt3}\left(\frac{137}{2\pi^2}\right)^{1/3},
 \qquad
 \|v\|_{L^6(L(137,1))}
 \le C_{L(137,1)}\|v\|_{H_A^1}
 \le C_V\|v\|_{H_A^1}
 \qquad (R_c=1).
\end{equation}
This graph-form embedding is one-way: small $L^6$ norm does not imply small
graph norm, so an independently supplied graph-norm threshold $\delta$ cannot
be replaced by $\delta/C_V$ as an $L^6$ threshold.  The energy class and the Sobolev
embedding $V\hookrightarrow L^6(M)$ give,
for $0\le\theta\le1$,
\begin{equation}\label{eq:energy-interpolation}
 u\in L^p_{\mathrm{loc}}(L^q),\qquad
 \frac1p=\frac\theta2,\qquad
 \frac1q=\frac{1-\theta}{2}+\frac\theta6.
\end{equation}
Thus $2\le p\le\infty$ and $2\le q\le6$, with the exponents coupled by
\[
 \frac2p+\frac3q=\frac32,
\]
not the Serrin region $2/p+3/q\le1$ (with $q>3$ in its usual nonendpoint
form).  In particular, $(p,q)=(\infty,6)$ is not supplied by
\eqref{eq:energy-interpolation}: $p=\infty$ corresponds to $q=2$.  The exact
spectrum is infinite, so finite cutoff counts do not imply cascade starvation.

For the exact time-average statement this graph norm is equivalent to the
usual $H^1$ norm on the solenoidal sector, but the
constant below is asserted only for this declared graph norm.  Put
$E(t)=\|u(t)\|_H^2$ and $E_0=E(0)$, with no factor $1/2$.  For every
$t\in\mathcal G_u$ and $\Delta t>0$, the energy inequality gives an admissible
$s\in\mathcal G_u\cap[t,t+\Delta t]$ such that
\begin{equation}\label{eq:time-average}
 \|u(s)\|_{H_A^1}^2
 \le E(t)\left(1+\frac{1}{2\nu\Delta t}\right).
\end{equation}

The regularity consequence is explicitly conditional.  Assume an independent
theorem supplies $\delta>0$ for equation~\eqref{eq:hodge-ns}, the given
$(\nu,R_c)$, and this exact graph norm, such that every
$v\in V$ with $\|v\|_{H_A^1}\le\delta$ generates a global strong
continuation, and that weak--strong uniqueness identifies any Leray--Hopf
trajectory having datum $v$ with that continuation.  Then the original
trajectory is strong from the selected $s$ onward whenever the right-hand
side of~\eqref{eq:time-average} is at most $\delta^2$.  Since
$E(t)\le E_0e^{-2\nu\lambda_1t}$, this holds for every admissible $t$ satisfying
\begin{equation}\label{eq:conditional-onset}
 t\ge \frac{1}{2\nu\lambda_1}
 \log\!\left[
 \frac{E_0}{\delta^2}
 \left(1+\frac{1}{2\nu\Delta t}\right)
 \right],\qquad \lambda_1=\frac4{R_c^2}.
\end{equation}
The conditional onset is then no later than $t+\Delta t$.  Because
$\mathcal G_u$ has full measure, admissible $t$ may be chosen arbitrarily
close above the displayed threshold.  This is neither an unconditional
$T^*$ formula nor a theorem of global three-dimensional regularity or
uniqueness.

\section{Conclusion}\label{sec:conclusion}

We have combined the exact coexact spectrum of $L(137,1)$ with its
deck-restricted interaction algebra.  The result includes a
convention-complete Wigner formula, a discrete ladder-canonical phase, explicit
nonzero witnesses, and finite two-band data.  The same spectral gap yields
exponential energy decay and singleton phase and trajectory attractors for the
declared Hodge--Stokes evolution, without assuming uniqueness of weak
solutions.

Several questions remain open.  The finite tower census does not determine an
asymptotic cascade law, the canonical phase is not invariant under arbitrary
rephasing, and the conditional continuation argument does not prove global
three-dimensional regularity.  It would be particularly useful to understand
how much of the interaction algebra survives beyond round quotients and
whether the exact selection rules constrain long-range transfer across the
full, infinite spectrum.

\section*{Acknowledgment of assistance}
AI-assisted research, computation, checking, and editorial support used
OpenAI Codex and Anthropic Claude. Luqman Mahmood directed the work
and is responsible for the claims and final text. These internal
checks are not external peer review.

\begingroup
\footnotesize
\bibliographystyle{plain}
\bibliography{references}
\endgroup

\end{document}
